% chapters/12/praktikum.tex
\begin{praktikumbox}
Praktikum ini mengikuti Tugas Praktikum 09 pada RPS: memasang Unity LTS,
membuat \istilah{project} 2D, menyusun \istilah{scene} kosong berisi kamera dan
objek primitif, lalu menyimpannya.

\begin{langkah}
  \item \textbf{Pasang Unity.} Pasang Unity Hub, lalu pasang satu versi Unity
        \istilah{LTS} beserta modul dukungan 2D. Catat nomor versi yang Anda
        pasang pada lembar kerja, sebab nomor itu akan Anda pakai sepanjang
        blok ini.

  \item \textbf{Buat \istilah{project} 2D.} Melalui Unity Hub, buat
        \istilah{project} baru dengan templat 2D. Amati susunan panel yang
        muncul, lalu tentukan panel mana yang \istilah{Hierarchy}, mana yang
        \istilah{Scene}, mana yang \istilah{Inspector}, dan mana yang
        \istilah{Project window}.

  \item \textbf{Kenali scene bawaannya.} Periksa ketiga benda yang telah ada
        pada \istilah{scene} bawaan, lalu tentukan mana kamera, mana
        pencahayaan, dan mana objek primitif. Ubah letak objek primitif itu
        dengan menyeretnya pada panel \istilah{Scene}, lalu perhatikan perubahan
        nilai pada \istilah{Inspector}.

  \item \textbf{Tambahkan objek primitif.} Tambahkan tiga objek primitif yang
        berbeda, yaitu kubus, bola, dan silinder. Susun ketiganya berjajar, lalu
        atur warnanya agar berbeda-beda.

  \item \textbf{Tulis dan tempelkan skrip.} Ketik ulang
        \code{ch12_skrip_pertama.cs} dari Daftar~\ref{lst:ch12-skrip} ke dalam
        folder \code{Scripts} pada \istilah{Project window}, lalu tempelkan pada
        salah satu objek primitif.

  \item \textbf{Jalankan scene-nya.} Tekan tombol jalankan, lalu gerakkan objek
        itu dengan tombol panah. Amati pesan yang muncul pada panel
        \istilah{Console}. Catat selang bingkai dan jarak yang dilaporkan.

  \item \textbf{Simpan dan periksa berkasnya.} Simpan \istilah{scene} dan
        \istilah{project}, lalu periksa folder \code{Assets} pada penjelajah
        berkas komputer Anda. Catat berkas apa saja yang muncul di sana, dan
        berkas apa saja yang sebaiknya tidak disimpan pada repositori.

  \item \textbf{Tulis catatan praktikum} yang memuat nomor versi Unity,
        keterangan kegunaan setiap panel, tangkapan layar scene buatan Anda,
        dan tabel padanan konsep Pygame dan Unity yang Anda susun sendiri.

  \item \textbf{Simpan pekerjaan} ke repositori Git, dengan berkas
        \code{.gitignore} yang mendaftarkan folder \code{Library},
        \code{Temp}, \code{Obj}, dan \code{Build}.
\end{langkah}
\end{praktikumbox}

\begin{table}[htbp]
\centering
\caption{Lembar kerja Praktikum 09}
\label{tab:lembar-praktikum-09}
{\setlength{\tabcolsep}{4pt}\footnotesize%
\begin{tabularx}{\textwidth}{@{}L{4.2cm} Y@{}}
\toprule
\headrow \textbf{Butir yang dicatat} & \textbf{Keterangan} \\
\midrule
Nomor versi Unity yang dipasang & \\[1.2em]
\addlinespace
Kegunaan setiap panel, dengan kata-kata sendiri & \\[2.4em]
\addlinespace
Nama berkas yang muncul pada folder \code{Assets} & \\[1.2em]
\addlinespace
Selang bingkai yang dilaporkan pada panel \istilah{Console} & \\[1.2em]
\addlinespace
Kendala yang dihadapi beserta cara mengatasinya & \\[2.4em]
\bottomrule
\end{tabularx}}
\end{table}

\begin{catatabox}
Langkah kedelapan menuntut Anda menyusun tabel padanan dengan kata-kata
sendiri. Tabel itu bukan salinan dari Tabel~\ref{tab:padanan-lengkap}, melainkan
hasil pemeriksaan Anda sendiri: bukalah kembali kode Python Anda pada blok
sebelumnya, lalu tentukan padanannya pada Unity. Tabel itulah yang menjadi bukti
indikator \ind{c}.
\end{catatabox}
